\chapter{HASIL DAN PEMBAHASAN}

\section{Deskripsi Data dan Lingkungan Eksperimental}
\label{sec:deskripsi-data-lingkungan-eksperimental}

Eksperimen dilakukan untuk membandingkan metode representasi langkah laten
dan medium komunikasi pada \textit{Multi-Agent System} (MAS). Seluruh eksperimen
menggunakan model \textit{open-weight} Qwen3-8B sebagai \textit{backbone}, dengan
$L=36$ lapisan \textit{Transformer} \textit{decoder-only} dan dimensi
\textit{hidden state} $d=4096$. Model menggunakan \textit{grouped-query attention}
(GQA) dengan delapan \textit{key-value heads} serta \textit{untied embedding}.
Konfigurasi model tersebut dijaga tetap pada seluruh kondisi eksperimen.

Rancangan eksperimen menggunakan empat agen yang dijalankan secara sekuensial,
yaitu \textit{planner}, \textit{critic}, \textit{refiner}, dan \textit{judger}.
Pada medium \textit{KV-cache}, informasi internal diteruskan melalui
\textit{KV-cache}, sedangkan pada medium teks informasi diteruskan dalam bentuk
keluaran tekstual. Kondisi laten menggunakan hingga $m=10$ langkah laten dengan
persamaan $\Phi$ yang sesuai.

Eksperimen utama dilaksanakan pada 10 Agustus 2026 menggunakan tiga
\textit{benchmark}, yaitu GSM8K, ARC-Challenge, dan HumanEval+. GSM8K digunakan
untuk penalaran matematika, ARC-Challenge untuk tugas pilihan ganda, dan
HumanEval+ untuk menghasilkan program Python. Data dan tugas yang digunakan
disajikan pada Tabel~\ref{tab:data-benchmark}.

\begin{table}[H]
    \centering
    \caption{Data dan tugas yang digunakan dalam eksperimen utama.}
    \label{tab:data-benchmark}
    \small
    \begin{tabular}{llllr}
        \hline
        Tugas & Keluarga & Sumber & \textit{Split} \\
        \hline
        GSM8K & Matematika & \code{openai/gsm8k} & \code{test}\\
        ARC-C & Pilihan ganda & \code{allenai/ai2\_arc} & \code{test}\\
        HumanEval+ & Kode & \code{evalplus/humanevalplus} & \code{test}\\
        \hline
    \end{tabular}
\end{table}

Pada ARC-Challenge, label jawaban dikonversikan menjadi
huruf $a$ hingga $d$. Pada HumanEval+, keluaran program Python dievaluasi
menggunakan \textit{test harness} sesuai dengan \textit{entry point} pada
masing-masing soal.

Perlakuan utama terdiri atas empat metode langkah laten pada medium
\textit{KV-cache}, yaitu \code{raw}, \code{soft}, \code{sample}, dan
\code{moi}. Medium teks (\code{text}) digunakan sebagai kondisi pembanding.
Matriks perlakuan utama ditunjukkan pada Tabel~\ref{tab:matriks-perlakuan}.

\begin{table}[H]
    \centering
    \caption{Matriks perlakuan utama pada tiga \textit{benchmark}.}
    \label{tab:matriks-perlakuan}
    \small
    \begin{tabular}{lccccc}
        \hline
        Perlakuan & Medium & persamaan $\Phi$ & GSM8K & ARC-C & HumanEval+ \\
        \hline
        \code{raw}    & \code{kv}   & \code{raw}    & $\checkmark$ & $\checkmark$ & $\checkmark$ \\
        \code{soft}   & \code{kv}   & \code{soft}   & $\checkmark$ & $\checkmark$ & $\checkmark$ \\
        \code{sample} & \code{kv}   & \code{sample} & $\checkmark$ & $\checkmark$ & $\checkmark$ \\
        \code{moi}    & \code{kv}   & \code{moi}    & $\checkmark$ & $\checkmark$ & $\checkmark$ \\
        \code{text}   & \code{text} & ---           & $\checkmark$ & $\checkmark$ & $\checkmark$ \\
        \hline
    \end{tabular}
\end{table}

Parameter yang tidak menjadi objek variasi dijaga tetap selama eksperimen.
Ringkasan konfigurasi tersebut disajikan pada
Tabel~\ref{tab:konfigurasi-eksperimen}.

\begin{table}[H]
    \centering
    \caption{Konfigurasi tetap eksperimen utama.}
    \label{tab:konfigurasi-eksperimen}
    \small
    \begin{tabular}{ll}
        \hline
        Komponen & Nilai \\
        \hline
        \textit{Backbone} & Qwen3-8B \\
        Rantai agen & \textit{planner} $\rightarrow$ \textit{critic}
        $\rightarrow$ \textit{refiner} $\rightarrow$ \textit{judger} \\
        Jumlah langkah laten & $m=10$ \\
        Suhu langkah laten & $T=0{,}7$ \\
        Parameter MoI & $\beta=1$ \\
        Kriteria \textit{early stopping} &
        $\cos(h_{k+1},h_k)>0{,}999$ \\
        Suhu menghasilkan teks & $T_{\mathrm{gen}}=0{,}6$ \\
        \textit{Top-p} & $0{,}95$ \\
        Maksimum token keluaran & $2048$ \\
        Batas konteks medium \code{text} & $8000$ karakter \\
        \hline
    \end{tabular}
\end{table}

Dengan demikian, terdapat 15 sel eksperimen, yaitu lima kondisi perlakuan pada
tiga \textit{benchmark}, dengan 100 soal pada setiap sel. Seluruh kondisi dalam
\textit{benchmark} yang sama menggunakan himpunan soal identik. Ukuran yang
dicatat untuk setiap soal meliputi akurasi, \textit{format rate}, jumlah token
keluaran, dan waktu eksekusi.

%=============================
\section{Akurasi dan Keandalan Format}
\label{sec:akurasi-keandalan-format}

Evaluasi awal dilakukan untuk membandingkan akurasi dan \textit{correct format rate}
pada medium teks dan medium \textit{KV-cache}. Hasil perbandingan disajikan
pada Gambar~\ref{fig:b4-hasil-utama} dan Tabel~\ref{tab:akurasi-format}.

\begin{figure}[H]
    \centering
    \includegraphics[width=\textwidth]{build/assets/images/b4_g1_hasil_utama.png}
    \caption{Perbandingan akurasi, jumlah token keluaran per soal, dan waktu
    inferensi per soal pada tiga \textit{benchmark}.}
    \label{fig:b4-hasil-utama}
\end{figure}

Gambar~\ref{fig:b4-hasil-utama} menunjukkan bahwa akurasi metode
\textit{KV-cache} berada pada tingkat yang relatif dekat dengan medium teks
pada ketiga \textit{benchmark}. Pada GSM8K, medium teks menghasilkan akurasi
sebesar $0{,}90$, sedangkan \texttt{soft}, \texttt{sample}, dan \texttt{moi}
masing-masing menghasilkan $0{,}88$, $0{,}86$, dan $0{,}92$. Pada ARC-Challenge,
akurasi medium teks sebesar $0{,}93$, sementara \texttt{soft} dan
\texttt{sample} mencapai $0{,}94$ dan \texttt{moi} mencapai $0{,}91$.
Hasil tersebut menunjukkan bahwa penggunaan medium \textit{KV-cache} tidak
secara umum diikuti oleh penurunan akurasi yang besar.

Perbedaan lebih terlihat pada HumanEval+, di mana akurasi medium teks sebesar
$0{,}76$, sedangkan \texttt{soft}, \texttt{sample}, dan \texttt{moi}
masing-masing menghasilkan $0{,}69$, $0{,}74$, dan $0{,}72$. Dengan demikian,
penurunan akurasi juga terjadi pada metode laten lainnya, tetapi
\texttt{raw} menunjukkan penurunan yang lebih besar dengan akurasi sebesar
$0{,}42$.

Selain akurasi, \textit{correct format rate} digunakan untuk melihat keterpenuhan format
keluaran yang diperlukan dalam proses evaluasi. Hasilnya disajikan pada
Tabel~\ref{tab:akurasi-format}.

\begin{table}[H]
    \centering
    \caption{Akurasi dan \textit{correct format rate} pada masing-masing
    \textit{benchmark}.}
    \label{tab:akurasi-format}
    \small
    \begin{tabular}{llrrrrr}
        \hline
        Tolok ukur & Ukuran & Teks & \texttt{raw} & \texttt{soft}
        & \texttt{sample} & \texttt{moi} \\
        \hline
        GSM8K & Akurasi
        & 0,90 & 0,83 & 0,88 & 0,86 & \textbf{0,92} \\
        & \textit{correct format rate}
        & 1,00 & 0,98 & 1,00 & 0,97 & 1,00 \\
        \hline
        ARC-C & Akurasi
        & 0,93 & 0,89 & \textbf{0,94} & \textbf{0,94} & 0,91 \\
        & \textit{correct format rate}
        & 1,00 & 1,00 & 1,00 & 1,00 & 0,98 \\
        \hline
        HumanEval+ & Akurasi
        & 0,76 & 0,42 & 0,69 & \textbf{0,74} & 0,72 \\
        & \textit{correct format rate}
        & 1,00 & 1,00 & 0,96 & 1,00 & 1,00 \\
        \hline
    \end{tabular}
\end{table}

Berdasarkan Tabel~\ref{tab:akurasi-format}, \textit{correct format rate} berada pada
tingkat tinggi pada seluruh kondisi, dengan nilai antara $0{,}96$ dan
$1{,}00$. Hal ini menunjukkan bahwa perbedaan akurasi yang diamati tidak
secara umum disertai oleh kegagalan dalam memenuhi format keluaran. Pada
HumanEval+, misalnya, \texttt{raw} tetap memiliki \textit{correct format rate} sebesar
$1{,}00$ meskipun akurasinya hanya sebesar $0{,}42$.

%=============================
%=============================
\section{Biaya Token dan Waktu Inferensi}
\label{sec:biaya-token-waktu-inferensi}

Selain mempertahankan akurasi, penggunaan medium \textit{KV-cache} juga
ditinjau dari biaya inferensinya. Jumlah token keluaran digunakan untuk
menggambarkan banyaknya keluaran tekstual yang dihasilkan selama proses
komunikasi, sedangkan waktu per soal menunjukkan durasi yang diperlukan hingga
suatu soal selesai diproses. Perbandingan kedua ukuran tersebut ditunjukkan
pada Gambar~\ref{fig:b4-hasil-utama} dan Tabel~\ref{tab:biaya-inferensi}.

\begin{table}[H]
    \centering
    \caption{Jumlah token keluaran dan waktu inferensi per soal.}
    \label{tab:biaya-inferensi}
    \small
    \begin{tabular}{llrrrrr}
        \hline
        Tolok ukur & Ukuran & Teks & \texttt{raw} & \texttt{soft}
        & \texttt{sample} & \texttt{moi} \\
        \hline
        GSM8K & Token/soal
        & 844,7 & 211,5 & 192,3 & 187,4 & 189,8 \\
        & Detik/soal
        & 49,2 & 19,7 & 14,5 & 16,5 & 15,9 \\
        \hline
        ARC-C & Token/soal
        & 1081,0 & 219,4 & 154,1 & 145,8 & 145,0 \\
        & Detik/soal
        & 76,0 & 20,6 & 14,0 & 13,3 & 15,3 \\
        \hline
        HumanEval+ & Token/soal
        & 997,2 & 176,8 & 121,5 & 126,0 & 121,6 \\
        & Detik/soal
        & 69,3 & 17,8 & 11,4 & 12,4 & 14,1 \\
        \hline
    \end{tabular}
\end{table}

Berdasarkan Tabel~\ref{tab:biaya-inferensi}, medium \textit{KV-cache} menghasilkan
jumlah token keluaran yang lebih rendah dibandingkan medium teks pada seluruh
\textit{benchmark}. Pada GSM8K, misalnya, jumlah token per soal turun dari
$844{,}7$ pada medium teks menjadi sekitar $187{,}4$--$211{,}5$ token pada
medium \textit{KV-cache}. Pola serupa terlihat pada ARC-Challenge dan
HumanEval+, dengan seluruh metode \textit{KV-cache} menghasilkan keluaran
yang lebih sedikit dibandingkan medium teks.

Pengurangan jumlah token tersebut diikuti oleh penurunan waktu inferensi.
Pada GSM8K, waktu per soal berkurang dari $49{,}2$ detik pada medium teks
menjadi $14{,}5$--$19{,}7$ detik pada medium \textit{KV-cache}. Pada
ARC-Challenge dan HumanEval+, pola yang sama juga terlihat, yaitu seluruh
metode \textit{KV-cache} memerlukan waktu inferensi yang lebih rendah
daripada medium teks.

Pola tersebut menunjukkan bahwa perbedaan biaya inferensi terutama terlihat
pada tingkat medium komunikasi. Medium teks memerlukan pembuatan keluaran
tekstual yang lebih banyak, sedangkan \textit{KV-cache} memungkinkan informasi
antaragen diteruskan tanpa seluruh informasi tersebut dibangkitkan kembali
dalam bentuk teks. Meskipun terdapat perbedaan biaya di antara metode
\textit{KV-cache}, perbedaan tersebut relatif lebih kecil dibandingkan
perbedaan antara medium \textit{KV-cache} dan medium teks.

%=============================
%=============================
\section{Efisiensi Medium KV terhadap Medium Teks}
\label{sec:efisiensi-medium-kv-teks}

Perbandingan selanjutnya dilakukan untuk mengukur efisiensi medium
\textit{KV-cache} relatif terhadap medium teks. Efisiensi dinilai berdasarkan
penghematan jumlah \textit{token} keluaran dan percepatan waktu inferensi.
Kedua ukuran dihitung terhadap kondisi teks pada \textit{benchmark} yang sama.
Hasil perbandingan disajikan pada Tabel~\ref{tab:efisiensi-kv-teks} dan
Gambar~\ref{fig:b4-pareto}.

\begin{table}[H]
    \centering
    \caption{Efisiensi medium \textit{KV-cache} relatif terhadap medium teks.}
    \label{tab:efisiensi-kv-teks}
    \small
    \begin{tabular}{llrrr}
        \hline
        Tolok ukur & Langkah laten & Akurasi &
        Penghematan \textit{token} & Percepatan \\
        \hline
        GSM8K & \texttt{raw}    & 0,83 & 75,0\% & $\times 2,5$ \\
              & \texttt{soft}   & 0,88 & 77,2\% & $\times 3,4$ \\
              & \texttt{sample} & 0,86 & 77,8\% & $\times 3,0$ \\
              & \texttt{moi}    & 0,92 & 77,5\% & $\times 3,1$ \\
        \hline
        ARC-C & \texttt{raw}    & 0,89 & 79,7\% & $\times 3,7$ \\
              & \texttt{soft}   & 0,94 & 85,7\% & $\times 5,4$ \\
              & \texttt{sample} & 0,94 & 86,5\% & $\times 5,7$ \\
              & \texttt{moi}    & 0,91 & 86,6\% & $\times 5,0$ \\
        \hline
        HumanEval+ & \texttt{raw}    & 0,42 & 82,3\% & $\times 3,9$ \\
                   & \texttt{soft}   & 0,69 & 87,8\% & $\times 6,1$ \\
                   & \texttt{sample} & 0,74 & 87,4\% & $\times 5,6$ \\
                   & \texttt{moi}    & 0,72 & 87,8\% & $\times 4,9$ \\
        \hline
    \end{tabular}
\end{table}

Berdasarkan Tabel~\ref{tab:efisiensi-kv-teks}, seluruh metode
\textit{KV-cache} menghasilkan penghematan \textit{token} dan percepatan
waktu inferensi pada ketiga \textit{benchmark}. Penghematan \textit{token}
berada pada rentang $75{,}0\%$--$87{,}8\%$, sedangkan percepatan berada pada
rentang $\times 2{,}5$--$\times 6{,}1$. Dengan demikian, efisiensi medium
\textit{KV-cache} terlihat secara konsisten ketika dibandingkan dengan medium
teks, meskipun besarnya penghematan berbeda menurut \textit{benchmark} dan
metode yang digunakan.

Gambar~\ref{fig:b4-pareto} memperlihatkan hubungan antara akurasi dengan
jumlah \textit{token} dan waktu inferensi. Pada kedua panel, kondisi teks
berada pada wilayah dengan biaya inferensi yang lebih tinggi, sedangkan
metode \textit{KV-cache} berada pada wilayah dengan jumlah \textit{token}
dan waktu yang lebih rendah. Pola tersebut menunjukkan bahwa penggunaan
\textit{KV-cache} memberikan pengurangan biaya komunikasi sekaligus
mempertahankan tingkat akurasi yang telah dibahas pada bagian sebelumnya.

\begin{figure}[H]
    \centering
    \includegraphics[width=\textwidth]{build/assets/images/b4_g2_pareto.png}
    \caption{Hubungan akurasi dengan jumlah \textit{token} keluaran per soal
    dan waktu inferensi per soal.}
    \label{fig:b4-pareto}
\end{figure}

Efisiensi tersebut tidak sepenuhnya seragam pada seluruh metode. Pada
HumanEval+, misalnya, \texttt{raw} memiliki akurasi $0{,}42$ meskipun
menghasilkan penghematan \textit{token} sebesar $82{,}3\%$ dan percepatan
sebesar $\times 3{,}9$. Sementara itu, \texttt{moi} mencapai akurasi $0{,}72$
dengan penghematan \textit{token} sebesar $87{,}8\%$ dan percepatan sebesar
$\times 4{,}9$. Hal ini menunjukkan bahwa besarnya efisiensi perlu dibaca
bersama dengan akurasi, karena pengurangan biaya inferensi tidak dengan
sendirinya menunjukkan kesetaraan kinerja.

%=============================
\section{Perbandingan Berpasangan MoI terhadap metode lain}
\label{sec:perbandingan-berpasangan-moi-acuan}

Setelah membandingkan efisiensi medium \textit{KV-cache} terhadap medium teks,
analisis selanjutnya dilakukan secara berpasangan untuk mengetahui apakah
perbedaan akurasi \texttt{moi} terhadap kondisi acuan didukung oleh data pada
tingkat soal. Perbandingan dilakukan menggunakan soal yang sama pada setiap
pasangan kondisi. Oleh karena itu, setiap soal membentuk satu pasangan
observasi yang memungkinkan perbedaan kinerja dihitung secara langsung.

Dalam analisis ini, \texttt{moi} dibandingkan dengan dua kondisi acuan, yaitu
\texttt{raw} pada medium \textit{KV-cache} dan medium teks. Besarnya perbedaan
akurasi dinyatakan sebagai $\Delta$ akurasi, sedangkan ketidakpastian estimasi
ditunjukkan melalui selang kepercayaan (\textit{SK}) $95\%$. Untuk mengetahui
apakah terdapat bukti statistik mengenai perbedaan akurasi, digunakan uji
McNemar eksak. Hasil perbandingan tersebut disajikan pada
Tabel~\ref{tab:selisih-moi-acuan}.

\begin{table}[H]
    \centering
    \caption{Selisih akurasi berpasangan \texttt{moi} terhadap kondisi acuan.}
    \label{tab:selisih-moi-acuan}
    \small
    \begin{tabular}{llrr}
        \hline
        Tolok ukur & Acuan & $\Delta$ akurasi & SK 95\% \\
        \hline
        GSM8K & \texttt{raw} (KV)
        & +0,09 & $[+0,01;\,+0,17]$ \\
        GSM8K & Teks
        & +0,02 & $[-0,04;\,+0,08]$ \\
        \hline
        ARC-C & \texttt{raw} (KV)
        & +0,02 & $[-0,03;\,+0,07]$ \\
        ARC-C & Teks
        & -0,02 & $[-0,07;\,+0,03]$ \\
        \hline
        HumanEval+ & \texttt{raw} (KV)
        & +0,30 & $[+0,20;\,+0,40]$ \\
        HumanEval+ & Teks
        & -0,04 & $[-0,13;\,+0,05]$ \\
        \hline
    \end{tabular}
\end{table}

Berdasarkan Tabel~\ref{tab:selisih-moi-acuan}, perbedaan \texttt{moi} terhadap
\texttt{raw} menunjukkan pola yang berbeda pada ketiga \textit{benchmark}.
Pada GSM8K, \texttt{moi} memiliki akurasi $0{,}09$ lebih tinggi daripada
\texttt{raw}, dengan selang kepercayaan $95\%$ sebesar
$[+0{,}01;\,+0{,}17]$ dan nilai-$p$ sebesar $0{,}0490$. Pada ARC-C, selisih
tersebut lebih kecil, yaitu $+0{,}02$, dengan selang kepercayaan
$[-0{,}03;\,+0{,}07]$ dan nilai-$p$ sebesar $0{,}6875$. Perbedaan terbesar
terlihat pada HumanEval+, dengan selisih akurasi sebesar $+0{,}30$,
selang kepercayaan $[+0{,}20;\,+0{,}40]$, dan nilai-$p<0{,}0001$.

Hasil tersebut menunjukkan bahwa keunggulan \texttt{moi} terhadap
\texttt{raw} paling jelas pada HumanEval+. Pada GSM8K terdapat perbedaan
positif yang relatif lebih kecil, sedangkan pada ARC-C selisihnya sangat
kecil dan interval kepercayaannya mencakup nilai nol. Dengan demikian,
perbedaan kinerja antara kedua metode tidak memiliki pola yang sama pada
seluruh jenis tugas.

Perbandingan terhadap medium teks menunjukkan hasil yang berbeda. Pada
GSM8K, \texttt{moi} memiliki akurasi $0{,}02$ lebih tinggi daripada teks,
sedangkan pada ARC-C dan HumanEval+ masing-masing lebih rendah sebesar
$0{,}02$ dan $0{,}04$. Ketiga selang kepercayaan masing-masing mencakup nilai
nol, dengan nilai-$p$ sebesar $0{,}7539$, $0{,}6875$, dan $0{,}5235$.
Dengan demikian, berdasarkan pengujian berpasangan, tidak terdapat bukti
statistik yang cukup mengenai perbedaan akurasi antara \texttt{moi} dan
medium teks pada ketiga \textit{benchmark} tersebut.

Secara keseluruhan, hasil perbandingan menunjukkan bahwa \texttt{moi}
mempertahankan akurasi yang tidak berbeda secara signifikan dari medium teks,
sementara memberikan keunggulan yang lebih nyata dibandingkan \texttt{raw}
pada HumanEval+. Hasil ini memperlihatkan bahwa metode \texttt{moi} tidak
hanya perlu dinilai berdasarkan akurasi, tetapi juga berdasarkan mekanisme
representasi informasi yang mendasari perbedaan kinerja tersebut.

%=============================
\section{Geometri Vektor Langkah Laten}
\label{sec:geometri-vektor-langkah-laten}

Hasil pada Subbab~\ref{sec:perbandingan-berpasangan-moi-acuan} menunjukkan
bahwa perbedaan kinerja \texttt{moi} terhadap \texttt{raw} lebih jelas pada
HumanEval+ dibandingkan pada GSM8K dan ARC-C. Analisis selanjutnya digunakan
untuk melihat apakah perbedaan tersebut disertai oleh perbedaan karakteristik
geometris representasi langkah laten. Fokus analisis adalah kedekatan vektor
langkah laten terhadap ruang \textit{input embedding} token.

Kedekatan tersebut diukur menggunakan kemiripan kosinus maksimum antara vektor
langkah laten dan seluruh vektor pada matriks \textit{input embedding}, yaitu

\begin{equation}
    s\bigl(\Phi(h)\bigr)
    =
    \max_i
    \cos\left(
    \Phi(h),W_{\mathrm{in}}[i]
    \right),
    \label{eq:geometri-cosine}
\end{equation}

di mana nilai yang semakin mendekati $1$ menunjukkan bahwa terdapat vektor
\textit{embedding} token dengan arah yang semakin dekat terhadap vektor langkah
laten. Oleh karena itu, nilai ini digunakan sebagai ukuran kedekatan geometris
representasi laten terhadap ruang representasi token.

\begin{table}[H]
    \centering
    \caption{Geometri vektor langkah laten berdasarkan kedekatan terhadap
    \textit{input embedding} token.}
    \label{tab:geometri-vektor}
    \small
    \begin{tabular}{l l c c c}
        \hline
        Langkah laten & Vektor yang disisipkan & Konveks &
        $\max_i\cos$ & Rentang \\
        \hline
        \texttt{raw}
        & $\displaystyle \tau\,\frac{hM}{\lVert hM\rVert}$
        & tidak & 0,312 & $[0,164;\,0,385]$ \\

        \texttt{soft}
        & $w=p$
        & ya & 0,927 & $[0,784;\,1,000]$ \\

        \texttt{sample}
        & $w=e_y$
        & ya & 1,000 & $[1,000;\,1,000]$ \\

        \texttt{moi}
        & $\displaystyle
        w=
        \frac{\tilde{H}p+(\beta+1-\tilde{H})e_y}
        {\beta+1}$
        & ya & 1,000 & $[1,000;\,1,000]$ \\
        \hline
    \end{tabular}
\end{table}

Tabel~\ref{tab:geometri-vektor} memperlihatkan pemisahan yang jelas antara
\texttt{raw} dan tiga metode lainnya. Nilai $\max_i\cos$ pada \texttt{raw}
hanya sebesar $0{,}312$ dengan rentang $[0{,}164;\,0{,}385]$. Sebaliknya,
\texttt{soft} memiliki nilai $0{,}927$, sedangkan \texttt{sample} dan
\texttt{moi} mencapai $1{,}000$ dengan rentang $[1{,}000;\,1{,}000]$.
Dengan demikian, representasi yang dihasilkan oleh \texttt{soft},
\texttt{sample}, dan \texttt{moi} memiliki kedekatan yang jauh lebih tinggi
terhadap \textit{embedding} token dibandingkan \texttt{raw}.

Perbedaan tersebut sejalan dengan struktur persamaan yang digunakan oleh
masing-masing metode. Pada \texttt{soft}, \texttt{sample}, dan \texttt{moi},
representasi dibentuk dari bobot yang berada pada ruang konveks
\textit{embedding} token. Sifat tersebut membatasi representasi yang dihasilkan
pada ruang yang dibentuk oleh vektor-vektor \textit{embedding} tersebut.
Khusus pada \texttt{sample}, representasi secara langsung berimpit dengan
vektor \textit{embedding} token, sehingga nilai kemiripan maksimum mencapai
$1{,}000$. Pada \texttt{moi}, nilai yang sama menunjukkan bahwa pada konfigurasi
yang diamati terdapat vektor \textit{embedding} yang memiliki arah berimpit
dengan representasi yang dihasilkan.

Sebaliknya, \texttt{raw} tidak memiliki batasan geometris yang sama. Nilai
kemiripan maksimum yang lebih rendah menunjukkan bahwa representasi
\texttt{raw} cenderung berada lebih jauh dari vektor \textit{embedding} token.
Pola tersebut juga terlihat pada Gambar~\ref{fig:b4-geometri}, yang
memperlihatkan pemisahan antara \texttt{raw} dan kelompok
\texttt{soft}, \texttt{sample}, dan \texttt{moi}.

\begin{figure}[H]
    \centering
    \includegraphics[width=\textwidth]{build/assets/images/b4_g3_geometri.png}
    \caption{Kedekatan geometris vektor langkah laten terhadap
    \textit{input embedding} token.}
    \label{fig:b4-geometri}
\end{figure}

Untuk menilai apakah rendahnya kedekatan \texttt{raw} dapat dijelaskan hanya
oleh kemiripan antara ruang sebelum dan sesudah transformasi, digunakan pula
beberapa ukuran diagnostik. Jarak relatif transformasi terhadap matriks
identitas diperoleh sebesar
\begin{equation}
    \frac{\lVert M-I\rVert_F}{\lVert I\rVert_F}
    =
    1{,}04,
    \label{eq:geometri-matrix-diagnostic}
\end{equation}

sedangkan rata-rata kemiripan arah antara $h$ dan $hM$ adalah

\begin{equation}
    \overline{\cos(h,hM)}
    =
    0{,}0112.
    \label{eq:geometri-h-hm}
\end{equation}

Nilai \(\lVert M-I\rVert_F/\lVert I\rVert_F=1{,}04\) menunjukkan bahwa beda antara 
matriks transformasi \(M\) dan matriks identitas \(I\) sedikit lebih besar daripada 
ukuran matriks identitas itu sendiri. Jadi, \(M\) tidak bisa dibilang dekat dengan 
identitas kalau dilihat dari ukuran ini. Hal ini juga terlihat dari rata-rata kemiripan 
arah antara \(h\) dan \(hM\) yang hanya \(0{,}0112\). Angka itu sangat dekat dengan nol, 
jadi \(hM\) secara rata-rata tidak lagi searah dengan \(h\). Kedua angka ini sama-sama 
menunjukkan bahwa transformasi pada \texttt{raw} benar-benar mengubah representasi data, 
bukan hanya mengubah skala atau mempertahankan arah vektor awalnya.

Selain itu, rata-rata kemiripan antara \textit{input embedding} dan
\textit{output embedding} pada indeks yang sama adalah

\begin{equation}
    \overline{
    \cos\left(
    W_{\mathrm{in}}[i],W_{\mathrm{out}}[i]
    \right)
    }
    =
    0{,}0041.
    \label{eq:geometri-in-out}
\end{equation}

Nilai $0{,}0041$ yang juga sangat dekat dengan nol menunjukkan bahwa
\textit{input embedding} dan \textit{output embedding} pada indeks yang sama
memiliki kesamaan arah yang sangat rendah. Ini penting karena menunjukkan bahwa
kedekatan vektor \texttt{raw} dengan \textit{input embedding} tidak bisa
dijelaskan hanya dengan anggapan bahwa kedua ruang \textit{embedding} itu
punya arah yang hampir sama. Jadi, ketiga ukuran pemeriksaan ini bersama-sama
menunjukkan bahwa transformasi dan hubungan antarruang \textit{embedding} pada
model yang dipakai memiliki perbedaan arah yang cukup besar.

Ketiga ukuran pemeriksaan ini menguatkan bahwa kedekatan \texttt{raw} yang
rendah dengan \textit{input embedding} bukan hanya akibat transformasi yang
mempertahankan arah vektor secara langsung.

Temuan ini menjadi konteks untuk hasil perbandingan berpasangan pada
Subbab~\ref{sec:perbandingan-berpasangan-moi-acuan}. Pada HumanEval+,
\texttt{moi} memiliki selisih akurasi $0{,}30$ terhadap \texttt{raw},
sedangkan perbedaannya dengan medium teks tidak cukup bukti secara statistik.
Pola ini cocok dengan kemungkinan bahwa perbedaan kinerja antara \texttt{moi}
dan \texttt{raw} berkaitan dengan perbedaan sifat geometris representasi laten
yang dihasilkan kedua metode.

Namun, kedekatan geometris saja tidak berarti ada hubungan sebab-akibat antara
representasi yang lebih dekat dengan \textit{embedding} token dan akurasi.
Hasil pada subbab ini menunjukkan adanya perbedaan sifat representasi
antarmetode, sedangkan hubungan antara perubahan geometri dan perubahan
akurasi perlu dilihat lewat variasi representasi yang dikontrol lebih khusus.

%=============================
\section{Kedekatan Geometris dan Akurasi}
\label{sec:kedekatan-geometris-akurasi}

\begin{figure}[htbp]
    \centering
    \includegraphics[width=0.95\textwidth]{build/assets/images/b4_g6_ambang.png}
    \caption{Hubungan kedekatan geometris dengan akurasi HumanEval+ pada
    sumbu interpolasi dari \code{raw} menuju \code{soft}.}
    \label{fig:ambang-geometri}
\end{figure}

Hasil pada Subbab~\ref{sec:geometri-vektor-langkah-laten} menunjukkan adanya
perbedaan kedekatan geometris antara vektor langkah laten dan
\textit{embedding} token, terutama antara \code{raw} dan metode lainnya.
Untuk melihat apakah perubahan kedekatan geometris berkaitan dengan perubahan
akurasi, dilakukan percobaan interpolasi dari representasi \code{raw} menuju
\code{soft}. Percobaan ini digunakan sebagai analisis tambahan untuk melihat
perubahan akurasi ketika representasi \code{raw} secara bertahap digeser
menuju representasi \code{soft}.

Hasil interpolasi disajikan pada Gambar~\ref{fig:ambang-geometri} dan
Tabel~\ref{tab:b4-ambang}. Nilai $\alpha=0$ menunjukkan representasi \code{raw},
sedangkan $\alpha=1$ menunjukkan representasi \code{soft}. Kedekatan
geometris diukur dengan nilai maksimum kemiripan kosinus terhadap seluruh
\textit{input embedding} token. Semakin besar nilainya, semakin dekat
representasi laten dengan arah salah satu \textit{embedding} token.

\begin{equation}
    s(\alpha)
    =
    \max_i
    \cos\left(
    z_{\alpha},W_{\mathrm{in}}[i]
    \right).
    \label{eq:kedekatan-geometris}
\end{equation}

\input{build/assets/tables/b4_ambang}

Berdasarkan Tabel~\ref{tab:b4-ambang}, pada $\alpha=0$, representasi
\code{raw} memiliki kedekatan geometris sebesar $0{,}312$ dan akurasi
HumanEval+ sebesar $0{,}42$. Ketika $\alpha$ dinaikkan menjadi $0{,}25$,
kedekatan geometris meningkat menjadi $0{,}452$, tetapi akurasi justru turun
menjadi $0{,}28$. Jadi, peningkatan kedekatan geometris pada tahap awal tidak
langsung diikuti oleh peningkatan akurasi.

Perubahan yang lebih jelas terlihat mulai dari $\alpha=0{,}50$. Pada nilai
tersebut, kedekatan geometris meningkat menjadi $0{,}720$ dan akurasi naik
menjadi $0{,}68$. Selanjutnya, pada $\alpha=0{,}75$, kedekatan geometris
mencapai $0{,}924$ dengan akurasi $0{,}69$. Pada $\alpha=1$, kedekatan
geometris hanya meningkat sedikit menjadi $0{,}927$, sedangkan akurasi tetap
sebesar $0{,}69$.

Pola tersebut menunjukkan bahwa perubahan akurasi tidak mengikuti kedekatan
geometris secara linear. Pada tahap awal, kenaikan kedekatan dari $0{,}312$
menjadi $0{,}452$ justru disertai penurunan akurasi. Sebaliknya, setelah
kedekatan mencapai $0{,}720$, akurasi meningkat dari $0{,}28$ menjadi
$0{,}68$ dan kemudian berada pada kisaran $0{,}68$--$0{,}69$. Pada tingkat
kedekatan sekitar $0{,}92$, peningkatan kedekatan lebih lanjut tidak lagi
diikuti perubahan akurasi yang terlihat.

Dengan demikian, hasil interpolasi menunjukkan adanya hubungan antara posisi
geometris representasi dan akurasi HumanEval+. Perubahan akurasi yang lebih besar muncul
ketika representasi bergerak cukup jauh dari \code{raw} menuju \code{soft}.
Hal ini mendukung dugaan bahwa ketidakselarasan geometris \code{raw} dengan
ruang \textit{embedding} token berkaitan dengan kinerjanya pada HumanEval+,
meskipun hasil ini belum cukup untuk menyatakan bahwa kedekatan geometris
merupakan penyebab langsung perubahan akurasi.

%=============================
\section{Ketakpekaan MoI terhadap $\beta$}
\label{sec:ketakpekaan-moi-beta}

Parameter $\beta$ pada metode \code{moi} mengatur bobot relatif antara
distribusi $p$ dan token terpilih $e_y$. Berdasarkan
Persamaan~\eqref{eq:moi-beta-influence}, koefisien distribusi \textit{soft} dapat
dituliskan sebagai

\begin{equation}
a=\frac{\tilde{H}}{\beta+1},
\label{eq:beta-koefisien}
\end{equation}

sehingga peningkatan $\beta$ memperkecil kontribusi $p$ dan membuat
representasi semakin didominasi oleh $e_y$. Untuk melihat sensitivitas hasil
terhadap parameter tersebut, dilakukan variasi $\beta$ pada rentang
$0{,}25$ hingga $8{,}00$. Hasil pengujian disajikan pada
Tabel~\ref{tab:b4-beta} dan Gambar~\ref{fig:beta}.

\begin{figure}[htbp]
    \centering
    \includegraphics[width=0.95\textwidth]{build/assets/images/b4_g5_beta.png}
    \caption{Recall informasi simbolik pada berbagai nilai $\beta$.}
    \label{fig:beta}
\end{figure}

\input{build/assets/tables/b4_beta}

Berdasarkan Tabel~\ref{tab:b4-beta}, perubahan $\beta$ tidak menghasilkan
perubahan yang besar pada \textit{recall}. Pada \textit{recall} nama fungsi,
nilainya berada pada rentang $0{,}35$--$0{,}38$, sedangkan \textit{recall}
token acak berada pada rentang $0{,}11$--$0{,}13$. Dengan demikian, meskipun
$\beta$ mengubah bobot komponen penyusun representasi \code{moi}, perubahan
tersebut tidak memberikan pengaruh yang signifikan.

Hasil ini menunjukkan bahwa kinerja kanal laten \texttt{moi} tidak banyak berubah 
saat nilai $\beta$ diubah pada rentang yang diuji. Jadi, nilai $\beta=1$ sebagai 
nilai standar masih bisa dipakai, tanpa perlu mencari parameter khusus untuk 
mendapatkan kinerja kanal yang jelas berbeda.

%=============================
\section{Kinerja Metode Laten}
\label{sec:kinerja-metode-laten}

Selain akurasi dan efisiensi, fidelitas keluaran juga diperiksa melalui
kemunculan jawaban dengan aksara non-Latin. Hasil pada
Tabel~\ref{tab:b4-korupsi} menunjukkan bahwa \code{raw} menghasilkan enam
kejadian, sedangkan \code{sample} tidak menghasilkan kejadian tersebut.
Pada \code{soft} dan \code{moi}, masing-masing hanya terdapat satu kejadian.
Hal ini menunjukkan bahwa perbedaan metode tidak hanya terlihat dari nilai
akurasi dan efisiensi, tetapi juga pada bentuk keluaran yang dihasilkan.

\input{build/assets/tables/b4_korupsi}

Untuk melihat kinerja setiap metode secara keseluruhan, hasil akurasi,
jumlah \textit{token}, waktu inferensi, dan fidelitas keluaran dirangkum pada
Tabel~\ref{tab:b4-ringkas}.

\input{build/assets/tables/b4_ringkas}

Berdasarkan Tabel~\ref{tab:b4-ringkas}, acuan teks menghasilkan akurasi rerata
sebesar $0{,}863$ dengan kebutuhan $974{,}3$ \textit{token} dan waktu
inferensi $64{,}8$ detik per soal. Sementara itu, \code{moi} menghasilkan
akurasi rerata $0{,}850$ dengan hanya $152{,}1$ \textit{token} dan waktu
inferensi $15{,}1$ detik per soal. Artinya, dibandingkan acuan teks,
\code{moi} mengurangi jumlah \textit{token} sekitar $84{,}4\%$ dan waktu
inferensi sekitar $76{,}7\%$, dengan akurasi rerata yang tetap dekat dengan
acuan.

Akurasi MoI relatif tertinggi jika dibandingkan dengan metode laten lainnya, $0{,}837$ pada \code{soft},
$0{,}847$ pada \code{sample}, dan $0{,}713$ pada \code{raw}. Selain itu,
\code{moi} menggunakan jumlah \textit{token} paling sedikit, yaitu $152{,}1$
per soal, serta hanya menghasilkan satu jawaban beraksara non-Latin.
Kombinasi tersebut menunjukkan bahwa \code{moi} mampu mempertahankan kinerja
mendekati acuan teks dengan kebutuhan sumber daya yang jauh lebih
rendah. Hasil ini juga melengkapi analisis sebelumnya mengenai kedekatan
geometris representasi dan kapasitas kanal laten, yang menunjukkan bahwa
karakteristik representasi \code{moi} berkaitan dengan kemampuannya dalam
mempertahankan informasi sekaligus mengurangi biaya komputasi.